\documentclass[aps,pra,twocolumn,superscriptaddress,nofootinbib]{revtex4-2}

\usepackage{amsmath,amssymb}
\usepackage{bm}
\usepackage{graphicx}
\usepackage{xcolor}
\usepackage[colorlinks=true,linkcolor=blue!60!black,citecolor=blue!60!black,urlcolor=blue!60!black]{hyperref}
\usepackage{cleveref}

% ---------------------------------------------------------------------
% Placeholders. \PH{...} marks text or a number still to be written or
% checked; search for "\PH" before submission.
% ---------------------------------------------------------------------
\newcommand{\PH}[1]{\textcolor{red}{[\textsf{#1}]}}
% A figure file that does not exist yet is replaced by a framed placeholder.
\newcommand{\figorph}[2][\columnwidth]{%
  \IfFileExists{#2}{\includegraphics[width=#1]{#2}}%
  {\fbox{\parbox[c][0.5#1][c]{0.95#1}{\centering\PH{figure #2}}}}}

% ---------------------------------------------------------------------
% Numbers. figures/numbers.tex is written by figures/plot_first_draft.py
% from the simulation data (make doc first_draft compute). Every number
% without data stays a red placeholder.
% ---------------------------------------------------------------------
\InputIfFileExists{figures/numbers.tex}{}{}
\newcommand{\num}[1]{\ifcsname Num#1\endcsname\csname Num#1\endcsname\else\PH{#1}\fi}

% ---------------------------------------------------------------------
% Notation (repository notation; one macro per symbol so that it can be
% renamed in one place at the notation-alignment stage).
% ---------------------------------------------------------------------
\newcommand{\nq}{n_q}                      % number of qubits
\newcommand{\tr}{t_r}                      % reversal time
\newcommand{\tscan}{t_{\rm scan}}          % kicks of the scan
\newcommand{\tE}{t_E}                      % Ehrenfest time
\newcommand{\thz}{\theta_0}                % initial polar angle
\newcommand{\phz}{\phi_0}                  % initial azimuth
\newcommand{\fa}{\vartheta}                % flip angle
\newcommand{\iflip}{i_{\rm flip}}          % flipped qubit
\newcommand{\PR}{P}                        % Husimi participation ratio
\newcommand{\Fl}{\overline{F}_{\rm late}}  % late-window mean fidelity
\newcommand{\Pe}{\overline{P}_{\rm early}} % early-window mean participation ratio
\newcommand{\Pl}{\overline{P}_{\rm late}}  % late-window mean participation ratio
\newcommand{\wS}{w_S}                      % weight in the symmetric subspace
\newcommand{\NQ}{\mathcal N_Q}             % husimi_norm
\newcommand{\Ub}{U_r}                      % time-reversal (backward) step
\newcommand{\Uc}{U}                        % forward step (spin j and chain)
\newcommand{\Rf}{R}                        % butterfly rotation

\begin{document}

\title{Butterfly echoes in the quantum kicked top: island versus chaotic-sea initial states\\
  \PH{working title}}

\author{\PH{authors}}
\affiliation{\PH{affiliations}}

\date{\today}

\begin{abstract}
We study the Loschmidt echo of the quantum kicked top, realized as a chain of $\nq$ qubits with an all-to-all Ising
coupling, when a single qubit is rotated between an exact forward and an exact backward evolution. The two initial
spin coherent states are not chosen by hand but selected by a scan over the whole sphere that combines the classical
Lyapunov exponent with the fidelity and the Husimi participation ratio of the quantum evolution: one lies in a stable
island, the other in the chaotic sea, and the latter is checked not to be trapped in a sticky region. \PH{one or two
sentences on the main result: contrast of the echo fidelity and of the Husimi snapshots between the two states; role
of the Ehrenfest time; leakage out of the symmetric subspace.}
\end{abstract}

\maketitle

%=====================================================================
\section{Introduction}\label{sec:intro}
%=====================================================================

Classical chaos is defined by the exponential sensitivity of trajectories to their initial conditions. Quantum
mechanics has no trajectories, and the overlap of two states evolved by the same unitary is conserved, so this
definition does not carry over~\PH{cite Peres; Haake}. Peres proposed instead to compare the evolution of one initial
state under two slightly different dynamics~\PH{cite Peres 1984}. Equivalently, one evolves the state forward, applies
a perturbation, and evolves backward: the overlap of the returned state with the initial one, the Loschmidt
echo or fidelity, measures how strongly the perturbation is amplified by the dynamics~\PH{cite Jalabert--Pastawski;
Gorin et al.\ review; Goussev et al.\ Scholarpedia}.

If the perturbation is a local operation applied at the turning point, the echo protocol is the ``butterfly'' version
of the Loschmidt echo, closely related to out-of-time-order correlators and to operator spreading in many-body
systems~\PH{cite Larkin--Ovchinnikov; Swingle; Hashimoto; Schmitt et al.}. It has been realized with nuclear spins,
trapped ions and superconducting qubits~\PH{cite experiments, incl.\ the reference paper ``Time Reversal of Complex
Evolution on a Quantum Computer''}. A controlled setting for such protocols is the quantum kicked top~\PH{cite Haake,
Ku\'s, Scharf 1987; Chaudhury et al.\ 2009}: a single large spin whose classical limit has a mixed phase space with
stable islands and a chaotic sea, and which can be written as a chain of qubits so that time reversal and local
perturbations are natural operations.

In a mixed phase space the outcome of the echo depends on where the initial state sits. This work contrasts two
initial spin coherent states, one in a stable island and one in the chaotic sea. Rather than choosing them by hand, we
select them by a scan over the whole sphere that combines classical and quantum criteria, and we verify that the
chaotic state is not trapped in a sticky region near an island. \PH{summary of results and outline of the paper.}

%=====================================================================
\section{Model}\label{sec:model}
%=====================================================================

\subsection{Classical map}\label{sec:classical}

The classical top is a unit vector
\begin{equation}\label{eq:s}
  \bm s = (X,Y,Z) = (\sin\theta\cos\phi,\ \sin\theta\sin\phi,\ \cos\theta)
\end{equation}
on the unit sphere, with polar angle $\theta\in[0,\pi]$ and azimuth $\phi\in(-\pi,\pi]$. The pair $(\phi,Z)$ is
canonical, so phase space is the rectangle $-\pi<\phi\le\pi$, $-1\le Z\le1$ of area $4\pi$. One period of the map
(\texttt{src/classical/map\_mod.f90}) is a kick, a rotation about the $x$ axis by the kick angle $h$,
\begin{equation}\label{eq:rot}
  X' = X,\ \ Y' = Y\cos h - Z\sin h,\ \ Z' = Y\sin h + Z\cos h,
\end{equation}
followed by a torsion, a rotation about $z$ by an angle proportional to the new height $Z'$,
\begin{equation}\label{eq:tor}
  \bm s_{\rm new} = R_z(kZ')\,\bm s',
\end{equation}
where $k$ is the torsion strength per period and $R_z(\alpha)$ the rotation about $z$ by $\alpha$. Time $t=0,1,2,\dots$
counts kicks. The map preserves area on the sphere.

The finite-time Lyapunov exponent $\Lambda$ of an initial point $\bm s_0$ is the mean exponential growth rate of a
tangent vector $\delta\bm s_t$ over $N_\Lambda$ kicks,
\begin{equation}\label{eq:lyap}
  \Lambda = \frac{1}{N_\Lambda}\ln\frac{|\delta\bm s_{N_\Lambda}|}{|\delta\bm s_0|},
\end{equation}
computed with the tangent map of \cref{eq:rot,eq:tor} and periodic renormalization (Benettin's method,
\texttt{src/classical/lyapunov\_mod.f90}). Points in a stable island have $\Lambda\to0$, points in the chaotic sea
$\Lambda>0$.

\subsection{Quantum kicked top}\label{sec:quantum}

The quantum top is a spin $\bm S=(S_x,S_y,S_z)$ with $\bm S^2=j(j+1)$, $\hbar=1$, where $j$ is the spin quantum
number. One period is the Floquet operator (\texttt{src/quantum/floquet\_mod.f90})
\begin{equation}\label{eq:U}
  \begin{gathered}
  U = U_H(\varepsilon)\,U_K(h),\qquad \psi(t) = U^t\psi(0),\\
  U_H(\varepsilon)=e^{-i\varepsilon S_z^2},\qquad U_K(h)=e^{-ihS_x},
  \end{gathered}
\end{equation}
where $U_K(h)$ is the kick and $U_H(\varepsilon)$ the free evolution between kicks, with the coupling phase
$\varepsilon$ related to the classical torsion by
\begin{equation}\label{eq:eps}
  k = 2j\varepsilon .
\end{equation}
With $\bm s=\bm S/j$, the effective Planck constant is $\hbar_{\rm eff}=1/j$ and \cref{eq:U} reduces to the map
\cref{eq:rot,eq:tor} for $j\to\infty$. The kick acts first and the torsion second in both, so the $x$-kick convention
is the same classically and quantum mechanically.

A phase-space point $(\theta,\phi)$ corresponds to the spin coherent state
$|\theta,\phi\rangle = e^{-i\phi S_z}e^{-i\theta S_y}|j,j\rangle$, a spot of angular width $\sim1/\sqrt j$. The
Husimi distribution of a state $\psi$ is
\begin{equation}\label{eq:Q}
  Q(\theta,\phi) = |\langle\theta,\phi|\psi\rangle|^2,\qquad
  \frac{2j+1}{4\pi}\int Q\,d\Omega = \langle\psi|\psi\rangle ,
\end{equation}
with $d\Omega = d\phi\,dZ$ the area element of the sphere; it is normalized to one for a normalized spin-$j$ state.

\subsection{Qubit chain}\label{sec:chain}

The spin $j$ is realized by $\nq=2j$ qubits with Pauli matrices $\sigma^a_i$, $a\in\{x,y,z\}$, $i=1,\dots,\nq$, as
the collective spin
\begin{equation}\label{eq:Jq}
  S_a = \frac12\sum_{i=1}^{\nq}\sigma^a_i ,\qquad j=\frac{\nq}{2}.
\end{equation}
The kick factorizes into single-qubit pulses, $U_K(h)=e^{-ihS_x}=\prod_i e^{-ih\sigma^x_i/2}$, so $h$ is the
rotation angle of every qubit about $x$ \PH{the TIRKI model of the reference paper writes
$U_K(h)=\exp(-ih\sum_i\sigma^x_i)$, i.e.\ with half this angle; decide which convention to state}. The torsion is the
time-reversible kicked Ising (TIRKI) form of the reference paper with all-to-all couplings $J_{il}$,
\begin{equation}\label{eq:Hzz}
  H_0 = \sum_{i,l=1}^{\nq}J_{il}\,\sigma^z_i\sigma^z_l ,\qquad J_{il}=\frac14
  \quad\Rightarrow\quad H_0=S_z^2 ,
\end{equation}
where the sum includes $i=l$ (a constant $\nq/4$). One period of the chain (\texttt{src/chain/chain\_floquet\_mod.f90};
$J_{il}$ is \texttt{V\_il} in the code) is
\begin{equation}\label{eq:Uc}
  \begin{gathered}
  U = U_H(\varepsilon)\,U_K(h) = e^{-i\varepsilon H_0}\prod_{i=1}^{\nq}e^{-i\frac h2\sigma^x_i},\\
  \varepsilon=\frac{k}{2j}=\frac{k}{\nq},
  \end{gathered}
\end{equation}
which equals $U$ of \cref{eq:U} on the symmetric (Dicke) subspace of dimension $2j+1=\nq+1$, i.e.\ the
subspace of states invariant under permutations of the qubits. The coherent state $|\theta,\phi\rangle$ is the
product state with every qubit polarized along $\bm s(\theta,\phi)$.

\paragraph{Time reversal.}
The eigenvalues of $H_0$ are $m^2$ with $m=-j,\dots,j$, which differ by integers when $J_{il}=1/4$. Hence
$U_H(2\pi-\varepsilon) = \gamma_0\,U_H(\varepsilon)^{\dagger}$ with a global phase $\gamma_0$, and the time-reversal
operator
\begin{equation}\label{eq:Ub}
  U_r = U_K(-h)\,U_H(2\pi-\varepsilon) = \gamma\,U^{\dagger}
\end{equation}
inverts one forward step up to a global phase $\gamma$, using only a coupling of the same sign ($\varepsilon\to
2\pi-\varepsilon$) and pulses of the opposite sign ($h\to-h$).

%=====================================================================
\section{Methods}\label{sec:methods}
%=====================================================================

\subsection{Scan-based selection of the initial conditions}\label{sec:scan}

\paragraph{Grid.}
The initial states $|\thz,\phz\rangle$ are taken on an equal-area grid of $N_\theta\times N_\phi$ cell centers in
$(Z_0=\cos\thz,\ \phz)$, uniform in $Z_0\in[-1,1]$ and $\phz\in(-\pi,\pi]$, which covers the sphere with equal
weights. Here $\thz$ and $\phz$ are the polar angle and azimuth of the initial polarization. We use
$N_\theta\times N_\phi=\num{NthetaScan}\times\num{NphiScan}$ ($\num{Ngrid}$ states).

\paragraph{Classical pre-filter.}
For every grid point the finite-time Lyapunov exponent \cref{eq:lyap} is computed over $N_\Lambda=\num{Nlyap}$
kicks. Points with $\Lambda<\Lambda_{\rm reg}=\num{LamRegular}$ are island candidates, points with
$\Lambda>\Lambda_{\rm ch}=\num{LamChaotic}$ are chaotic-sea candidates.

\paragraph{Quantum scan.}
Every grid state is evolved by \cref{eq:U} for $\tscan$ kicks, recording the fidelity
\begin{equation}\label{eq:F}
  F(t) = |\langle\psi_0|\psi(t)\rangle|^2,\qquad \psi_0 = |\thz,\phz\rangle,
\end{equation}
the probability of finding the initial state again, and the Husimi participation ratio
\begin{equation}\label{eq:PR}
  \PR(t) = \frac{(2j+1)\left(\int Q\,d\Omega\right)^2}{4\pi\int Q^2\,d\Omega},
\end{equation}
the number of Planck cells of area $4\pi/(2j+1)$ covered by $Q$. A coherent state has $\PR=(4j+1)/(2j+1)\approx2$,
a random state of the spin $j$ has $\PR\approx j+1$ on average, and a uniform $Q$ would have $\PR=2j+1$.
\PR\ is computed without a phase-space grid: $|\theta,\phi\rangle\otimes|\theta,\phi\rangle$ is the spin-$2j$
coherent state, so $\int Q^2\,d\Omega = 4\pi|\chi|^2/(4j+1)$ with $\chi$ the projection of $\psi\otimes\psi$ onto
spin $2j$ \PH{short derivation in an appendix?}. For a chaotic state the quantum--classical correspondence holds up
to the Ehrenfest time
\begin{equation}\label{eq:tE}
  \tE \simeq \frac{|\ln\hbar_{\rm eff}|}{2\Lambda} = \frac{\ln j}{2\Lambda},
\end{equation}
the time for the exponential stretching $e^{\Lambda t}$ to blow the initial width $1/\sqrt j$ up to the size of the
sphere. We take $\tscan=\num{Tscan}\gg\tE$ (for the selected chaotic state $\tscan/\tE=\num{SeaTscanOverTe}$), and
average over two late windows, $\tscan/2<t\le3\tscan/4$ (early) and $3\tscan/4<t\le\tscan$ (late), giving $\Fl$,
$\Pe$ and $\Pl$.

\paragraph{Selection.}
The island state is the island candidate with the smallest $\Pl$. The chaotic state is the chaotic-sea candidate with
the largest $\Pl$ among those whose spread has saturated,
\begin{equation}\label{eq:sat}
  \left|\Pl/\Pe-1\right| < \delta_{\rm sat},\qquad \delta_{\rm sat}=\num{SatTol}.
\end{equation}
A state trapped in a sticky region near an island boundary, or localized by quantum interference, spreads slowly and
fails \cref{eq:sat}. We then check that the island state has $F$ close to one and that the chaotic state has $F$ close
to the random-state value $1/(2j+1)$.

\subsection{Butterfly echo on the chain}\label{sec:echo}

Each selected state is prepared as the product coherent state $\psi_0$ of $\nq$ qubits and evolved on the chain
(\texttt{src/combined/chain\_run\_mod.f90}): $\tr$ forward steps $\Uc$, the reversal time $\tr$ counting kicks,
then the single-qubit rotation
\begin{equation}\label{eq:R}
  \Rf = e^{-i\fa\sigma^a_{\iflip}/2} = \cos\frac{\fa}{2} - i\sin\frac{\fa}{2}\,\sigma^a_{\iflip},
\end{equation}
of qubit $\iflip\in\{1,\dots,\nq\}$ about the axis $a\in\{x,y,z\}$ by the flip angle $\fa$, and finally $\tr$ backward
steps $\Ub$:
\begin{equation}\label{eq:protocol}
  \psi(t) = \begin{cases} \Uc^{\,t}\psi_0, & 0\le t\le\tr,\\
  \Ub^{\,t-\tr}\,\Rf\,\Uc^{\,\tr}\psi_0, & \tr<t\le2\tr. \end{cases}
\end{equation}
Without the rotation ($\fa=0$) $\psi(2\tr)=\gamma^{\tr}\psi_0$, so $F(2\tr)=1$. With it, \cref{eq:Ub,eq:R} give
\begin{equation}\label{eq:Fecho}
  F(2\tr) = |\langle\psi(\tr)|\Rf|\psi(\tr)\rangle|^2
  = \cos^2\frac{\fa}{2} + \sin^2\frac{\fa}{2}\,\langle\sigma^a_{\iflip}\rangle_{\tr}^2 ,
\end{equation}
where $\langle\cdot\rangle_{\tr}$ is the expectation value in $\psi(\tr)$. The echo therefore measures the length of
the Bloch vector of the rotated qubit along $a$ at the turning point, which shrinks as the qubit becomes entangled with
the rest of the chain.

We record $F(t)$ for $0\le t\le2\tr$ and the Husimi distribution \cref{eq:Q} of the projection of $\psi(t)$ onto the
symmetric subspace at $t=0,\ \tr/2,\ \tr,\ 3\tr/2,\ 2\tr$. The rotation of one qubit breaks the permutation symmetry,
so after it part of the state leaves the symmetric subspace. Its weight there,
\begin{equation}\label{eq:wS}
  \wS(t) = \sum_{m=-j}^{j}|\langle j,m|\psi(t)\rangle|^2 ,
\end{equation}
with $|j,m\rangle$ the Dicke states, equals the Husimi norm $\NQ=\frac{2j+1}{4\pi}\int Q\,d\Omega$ of the projected
state (\texttt{husimi\_norm} in the code), so $\NQ<1$ after the flip. Since $\Ub$ and $\Uc$ commute with qubit
permutations, $\wS$ is constant on the backward leg. For a symmetric $\psi(\tr)$ the projector $P_S$ onto the symmetric subspace gives
$P_S\sigma^a_{\iflip}\psi(\tr)=\nq^{-1}\sum_iP_S\sigma^a_i\psi(\tr)=(S_a/j)\,\psi(\tr)$, so
$P_S\Rf\psi(\tr)=[\cos(\fa/2)-i\sin(\fa/2)S_a/j]\,\psi(\tr)$ and
\begin{equation}\label{eq:wSpred}
  \wS(t>\tr) = \cos^2\frac{\fa}{2} + \sin^2\frac{\fa}{2}\,\frac{\langle S_a^2\rangle_{\tr}}{j^2}.
\end{equation}
Optionally, the flip angle is scanned over $\fa=0,\pi/8,\dots,\pi$ for $a=x,y,z$ at fixed $\tr$, and $\tr$ is scanned
at fixed $(\fa,a)$.

\paragraph{Parameters.}
Throughout, $h=\num{Hval}$, $k=\num{Kval}$, $\nq=\num{NqVal}$ ($j=\num{Jval}$, $\varepsilon=k/\nq=\num{EpsVal}$),
$\tr=\num{Tr}$, $\iflip=\num{Iflip}$, $a=\num{FlipAxis}$ and $\fa=\pi$ unless stated otherwise.
\PH{justify the parameter choice: at $h=\pi/2$, $k=2.4$ no grid state stays localized at this $j$; at $h=1.2$ the
islands hold a coherent state and the chaotic sea covers a fraction \num{ChaoticFraction} of the sphere.}

%=====================================================================
\section{Results}\label{sec:results}
%=====================================================================

\subsection{Selected initial conditions}\label{sec:res-scan}

\begin{figure*}
  \centering
  \figorph[0.9\textwidth]{figures/scan.pdf}
  \caption{Initial-condition scan of the kicked top with kick angle $h=\num{Hval}$ rad about $x$, torsion strength
  $k=\num{Kval}$, spin $j=\num{Jval}$ (coupling $\varepsilon=k/(2j)=\num{EpsVal}$), on an equal-area grid of
  $\num{Ngrid}$ initial coherent states $|\theta_0,\phi_0\rangle$, plotted at their initial azimuth $\phi_0$ (rad,
  horizontal) and height $Z_0=\cos\theta_0$ (vertical). (a)~Classical finite-time Lyapunov exponent $\Lambda$ (per
  kick) over $N_\Lambda=\num{Nlyap}$ kicks. (b)~Quantum fidelity $F(t)=|\langle\psi_0|\psi(t)\rangle|^2$ averaged
  over the late window $3t_{\rm scan}/4<t\le t_{\rm scan}$, $t_{\rm scan}=\num{Tscan}$ kicks. (c)~Husimi participation
  ratio $P=(2j+1)(\int Q)^2/(4\pi\int Q^2)$ averaged over the same window, divided by its maximum $2j+1$; a coherent
  state has $P/(2j+1)\approx\num{PRcoherent}/(2j+1)$, a random state $\approx(j+1)/(2j+1)$. (d)~Ratio of the late-window to
  the early-window ($t_{\rm scan}/2<t\le3t_{\rm scan}/4$) mean of $P$; values near one mean that the spread has
  saturated. Stars: selected island state (blue) and chaotic-sea state (red).}
  \label{fig:scan}
\end{figure*}

\Cref{fig:scan} shows the four scan quantities over the sphere. \PH{describe: regular regions with $\Lambda\approx0$
coincide with high $\Fl$ and small $\Pl$; the chaotic sea with $\Lambda>0$, $\Fl\approx1/(2j+1)$ and large $\Pl$;
boundary regions with $\Pl/\Pe$ away from one.} The island state is $(\thz,\phz)=(\num{IslandThetaZero},
\num{IslandPhiZero})$ ($Z_0=\num{IslandZzero}$), with $\Lambda=\num{IslandLambda}$, $\Fl=\num{IslandFlate}$ and
$\Pl=\num{IslandPRlate}$. The chaotic state is $(\thz,\phz)=(\num{SeaThetaZero},\num{SeaPhiZero})$
($Z_0=\num{SeaZzero}$), with $\Lambda=\num{SeaLambda}$, $\Fl=\num{SeaFlate}$ (random-state value
$1/(2j+1)=\num{Frandom}$) and $\Pl=\num{SeaPRlate}$ (random-state value $j+1=\num{PRrandom}$). Its spread has
saturated, $\Pl/\Pe=\num{SeaPRratio}$, and its Ehrenfest time \cref{eq:tE} is $\tE=\num{SeaTehrenfest}$ kicks.

\begin{figure}
  \centering
  \figorph{figures/section.pdf}
  \caption{Classical Poincar\'e phase portrait of the kicked top with $h=\num{Hval}$ rad and $k=\num{Kval}$, in the
  canonical coordinates azimuth $\phi$ (rad) and height $Z=\cos\theta$. Black dots: orbits of
  $\num{NthetaSection}\times\num{NphiSection}$ equal-area initial points, $\num{NkicksSection}$ kicks each. Stars: the selected initial points, island
  $(\theta_0,\phi_0)=(\num{IslandThetaZero},\num{IslandPhiZero})$ (blue) and chaotic sea
  $(\num{SeaThetaZero},\num{SeaPhiZero})$ (red), with their own orbits in the same colors.}
  \label{fig:section}
\end{figure}

\begin{figure*}
  \centering
  \figorph[0.9\textwidth]{figures/scan_series.pdf}
  \caption{Spin-$j$ evolution $\psi(t)=U^t|\theta_0,\phi_0\rangle$, $U=e^{-i\varepsilon S_z^2}e^{-ihS_x}$, of the two
  selected initial states for $h=\num{Hval}$ rad, $k=\num{Kval}$, $j=\num{Jval}$ ($\varepsilon=\num{EpsVal}$). Blue: island
  state $(\theta_0,\phi_0)=(\num{IslandThetaZero},\num{IslandPhiZero})$; red: chaotic-sea state
  $(\num{SeaThetaZero},\num{SeaPhiZero})$. (a)~Fidelity $F(t)=|\langle\psi_0|\psi(t)\rangle|^2$; dotted: random-state
  value $1/(2j+1)$. (b)~Husimi participation ratio $P(t)$ (number of Planck cells); dotted: random-state value $j+1$.
  Light and dark gray bands: early and late averaging windows. Dashed red line: Ehrenfest time
  $t_E=\ln j/(2\Lambda)=\num{SeaTehrenfest}$ kicks of the chaotic state. Time $t$ in kicks.}
  \label{fig:scan-series}
\end{figure*}

\Cref{fig:section} places both states on the classical phase portrait. \PH{describe: the island state sits inside an
elliptic island \PH{near the period-1 fixed point?}, the chaotic state well inside the chaotic sea, away from island
boundaries.} \Cref{fig:scan-series} shows their full time series. \PH{describe: island $F$ oscillates close to one and
$P$ stays near the coherent value; chaotic $F$ drops to $\sim1/(2j+1)$ after $\sim\tE$ and $P$ rises to a plateau near
$j+1$ within a few $\tE$, with no slow drift.}

\subsection{Butterfly echo}\label{sec:res-echo}

\begin{figure*}
  \centering
  \figorph[0.9\textwidth]{figures/echo_fidelity.pdf}
  \caption{Butterfly echo on the chain of $n_q=\num{NqVal}$ qubits ($j=n_q/2$), $h=\num{Hval}$ rad,
  $\varepsilon=k/n_q=\num{EpsVal}$ ($k=\num{Kval}$), $J_{il}=1/4$: fidelity $F(t)=|\langle\psi_0|\psi(t)\rangle|^2$ versus
  time $t$ (kicks). Forward evolution for $t_r=\num{Tr}$ kicks (dotted vertical line), rotation
  $e^{-i\vartheta\sigma^a_{i_{\rm flip}}/2}$ of qubit $i_{\rm flip}=\num{Iflip}$ about $a=\num{FlipAxis}$ by
  $\vartheta=\pi$, then $t_r$ backward kicks. Circles: with the rotation; dashed: without it ($\vartheta=0$). Blue:
  island state $(\theta_0,\phi_0)=(\num{IslandThetaZero},\num{IslandPhiZero})$; red: chaotic-sea state
  $(\num{SeaThetaZero},\num{SeaPhiZero})$. Horizontal dotted line: $1/(2j+1)$. (a)~Linear, (b)~logarithmic scale.}
  \label{fig:echo}
\end{figure*}

\begin{figure*}
  \centering
  \figorph[\textwidth]{figures/echo_husimi.pdf}
  \caption{Husimi distribution $Q(\theta,\phi)=|\langle\theta,\phi|\psi_S(t)\rangle|^2$ of the projection $\psi_S$ of
  the chain state onto the symmetric subspace, at $t=0,\ t_r/2,\ t_r,\ 3t_r/2,\ 2t_r$ with $t_r=\num{Tr}$ kicks, for the
  echo of \cref{fig:echo} ($n_q=\num{NqVal}$, $h=\num{Hval}$ rad, $\varepsilon=\num{EpsVal}$, rotation of qubit
  $i_{\rm flip}=\num{Iflip}$ about $a=\num{FlipAxis}$ by $\vartheta=\pi$ at $t_r$). Top row: island state
  $(\theta_0,\phi_0)=(\num{IslandThetaZero},\num{IslandPhiZero})$; bottom row: chaotic-sea state
  $(\num{SeaThetaZero},\num{SeaPhiZero})$; white crosses mark $(\phi_0,Z_0)$. Horizontal axis: azimuth $\phi$ (rad);
  vertical axis: $Z=\cos\theta$; color: $Q$ (dimensionless), common scale within each row. After the rotation
  $(2j+1)\int Q\,d\Omega/(4\pi)=w_S<1$.}
  \label{fig:husimi}
\end{figure*}

\Cref{fig:echo} contrasts the two echoes. Without the rotation both return exactly, $|1-F(2\tr)|=\num{IslandFechoNoFlipDev}$ (island) and
$\num{SeaFechoNoFlipDev}$ (sea). With it, the island state has
$F(\tr)=\num{IslandFtr}$ and $F(2\tr)=\num{IslandFecho}$, and the chaotic state $F(\tr)=\num{SeaFtr}$ and
$F(2\tr)=\num{SeaFecho}$; both agree with \cref{eq:Fecho}, which gives \num{IslandFechoPred} and \num{SeaFechoPred}.
\PH{interpretation: revival versus decay; for $\fa=\pi$ the echo is $\langle\sigma^a_{\iflip}\rangle^2_{\tr}$, large
when the qubit stays polarized (island) and small when it is entangled (sea); the island echo depends on the
orientation of the polarization at $\tr$ relative to $a$.}

\Cref{fig:husimi} shows the Husimi snapshots. \PH{describe: island---$Q$ stays a compact spot that moves along the
island and returns to $(\phz,Z_0)$ at $2\tr$; sea---$Q$ covers the sphere at $\tr$, and after the rotation the
backward leg does not refocus it.} After the rotation the symmetric weight is $\wS=\num{IslandSymWeight}$ (island)
and $\num{SeaSymWeight}$ (sea), compared with \cref{eq:wSpred}: \num{IslandSymWeightPred} and \num{SeaSymWeightPred}.

\begin{figure*}
  \centering
  \figorph[0.9\textwidth]{figures/echo_scans.pdf}
  \caption{Echo fidelity $F(2t_r)$ on the chain ($n_q=\num{NqVal}$, $h=\num{Hval}$ rad, $\varepsilon=\num{EpsVal}$,
  rotated qubit $i_{\rm flip}=\num{Iflip}$). Blue: island state $(\theta_0,\phi_0)=(\num{IslandThetaZero},
  \num{IslandPhiZero})$; red: chaotic-sea state $(\num{SeaThetaZero},\num{SeaPhiZero})$. (a)~Versus flip angle
  $\vartheta$ (rad) for the axes $a=x$ (circles), $y$ (squares), $z$ (triangles) at $t_r=\num{Tr}$ kicks; lines:
  prediction $\cos^2(\vartheta/2)+\sin^2(\vartheta/2)\langle\sigma^a_{i_{\rm flip}}\rangle^2_{t_r}$. (b)~Versus reversal
  time $t_r$ (kicks) for $\vartheta=\pi$, $a=\num{FlipAxis}$: $F(2t_r)$ (circles, solid) and the weight $w_S(2t_r)$ of
  the symmetric subspace (squares, dashed); logarithmic vertical axis.}
  \label{fig:echo-scans}
\end{figure*}

\Cref{fig:echo-scans} \PH{describe the flip-angle/axis and $\tr$ dependence.}

%=====================================================================
\section{Discussion}\label{sec:discussion}
%=====================================================================

\paragraph{Lyapunov and Ehrenfest scaling.}
For the chaotic state the forward fidelity decays and the participation ratio saturates on the scale of the Ehrenfest
time \cref{eq:tE}, $\tE=\num{SeaTehrenfest}$ kicks for $\Lambda=\num{SeaLambda}$ and $j=\num{Jval}$. Since
$\tE\propto\ln j$, the window of quantum--classical correspondence grows only logarithmically with the number of qubits.
\PH{for $\tr$ below and above $\tE$: dependence of the echo on $\tr$ (\cref{fig:echo-scans}b); scaling with $\nq$ if
computed; relation of the decay of $\langle\sigma^a_{\iflip}\rangle_{\tr}^2$ to the single-qubit entanglement.} For the
island state $\Lambda\approx0$ and no such time scale exists: \PH{the echo is set by the regular motion of the
polarization}.

\paragraph{Symmetric-subspace leakage.}
The forward leg stays in the $(2j+1)$-dimensional symmetric subspace, where the chain is the kicked top. The rotation of
one qubit moves weight $1-\wS$ out of it, \cref{eq:wS,eq:wSpred}, and the Husimi norm of the projected state drops
below one, $\NQ=\num{IslandHusimiNormEnd}$ (island) and $\NQ=\num{SeaHusimiNormEnd}$ (sea) at $2\tr$. The backward
leg conserves this weight, so the leaked part never returns to the collective-spin description: \PH{interpretation;
which state leaks more and why ($\langle S_a^2\rangle_{\tr}/j^2$ close to one for a polarized state, smaller for a
spread state); implications for experiments that measure only collective observables.}

\paragraph{Outlook.}
\PH{larger $\nq$; static imperfections of amplitudes $\delta_x$ (kick) and $\delta_z$ (coupling) as in the reference paper; comparison
with the reference experiment; notation alignment with the reference paper.}

%=====================================================================
\section{Conclusion}\label{sec:conclusion}
%=====================================================================

\PH{conclusion}

\begin{acknowledgments}
\PH{acknowledgments}
\end{acknowledgments}

\PH{bibliography: \texttt{\textbackslash bibliography\{first\_draft\}} with a .bib file}

\end{document}
